\section{Utility Mechanisms}

\label{sec:utility}
% ============================================================================

\subsection{Compute Payment}
\label{sec:compute}

All AI workloads on Hanzo Network are priced in HANZO. The Hamiltonian Market
Maker (HMM)~\cite{hanzonetwork} clears resource markets by maintaining the
invariant:
\begin{equation}
\mathcal{H}(\bm{\Psi}, \bm{\Theta}) = \sum_{i=1}^{m} w_i \Psi_i \Theta_i +
\lambda \sum_{i=1}^{m} \tfrac{1}{2}(\Psi_i^2 + \Theta_i^2) = \kappa,
\end{equation}
where $\bm{\Psi}$ denotes compute supply (GPU-seconds, VRAM, bandwidth),
$\bm{\Theta}$ denotes demand credits, $w_i$ are resource weights, and $\lambda$
is a curvature parameter. Clients purchase demand credits with HANZO; workers
redeem credits for HANZO upon verified job completion.

\paragraph{Pricing.} The marginal price of resource $i$ is:
\begin{equation}
p_i = \frac{\partial \mathcal{H}}{\partial \Psi_i} \Big/
\frac{\partial \mathcal{H}}{\partial \Theta_i} =
\frac{w_i \Theta_i + \lambda \Psi_i}{w_i \Psi_i + \lambda \Theta_i}.
\end{equation}
This endogenous pricing eliminates dependence on external oracles.

\paragraph{Fee structure.} Each compute transaction incurs:
\begin{itemize}[leftmargin=1.1em]
    \item Market fee: 30 basis points (configurable by governance).
    \item Risk fee: 5--20 bps proportional to inventory imbalance.
    \item Of total fees collected, 25\% are burned and 75\% distributed to
    stakers and the treasury.
\end{itemize}

\subsection{Node Staking}
\label{sec:staking}

Running an inference node on Hanzo Network requires staking HANZO as collateral.
Staking serves three purposes: Sybil resistance, quality assurance (via
slashing), and supply reduction.

\begin{table}[H]
\centering
\small
\begin{tabular}{lrl}
\toprule
\textbf{Node Type} & \textbf{Min.\ Stake} & \textbf{Slashing} \\
\midrule
Validator & 100,000 HANZO & Up to 10\% \\
Worker (GPU) & Proportional to capacity & Up to 25\% \\
Verifier & 50,000 HANZO & Up to 5\% \\
Router & 10,000 HANZO & Up to 5\% \\
\bottomrule
\end{tabular}
\caption{Staking requirements by node type. Worker stake scales with declared GPU capacity.}
\label{tab:staking}
\end{table}

\paragraph{Staking rewards.} Node operators receive rewards from two sources:
\begin{enumerate}[leftmargin=1.4em]
    \item \textbf{Compute fees}: 50\% of platform fees (after burn) are
    distributed to stakers proportional to verified work.
    \item \textbf{Emission rewards}: From the community allocation, following the
    halving schedule described in \S\ref{sec:distribution}.
\end{enumerate}

\paragraph{Slashing conditions.} A staker's collateral is partially slashed if:
\begin{itemize}[leftmargin=1.1em]
    \item PoAI verification detects incorrect inference output.
    \item The node fails to respond within SLA timeout ($>$3 consecutive misses).
    \item The node serves a model version not approved by governance.
\end{itemize}

\subsection{Governance}
\label{sec:governance}

HANZO holders govern the network through an on-chain DAO with the following
scope:

\begin{itemize}[leftmargin=1.1em]
    \item \textbf{Model selection}: Which Zen models are served on the network
    (addition, deprecation, version upgrades).
    \item \textbf{Parameter tuning}: Fee rates, slashing percentages, staking
    minimums, emission schedule adjustments.
    \item \textbf{Treasury allocation}: Grant proposals, partnership funding,
    buyback programs.
    \item \textbf{Infrastructure upgrades}: Network upgrades, precompile
    activation, chain configuration changes.
\end{itemize}

\paragraph{Voting power.} Governance weight is computed as:
\begin{equation}
v_i = \alpha \cdot s_i + (1 - \alpha) \cdot h_i,
\end{equation}
where $s_i$ is the amount staked by participant $i$, $h_i$ is the amount held
(unstaked), and $\alpha = 0.7$ gives stakers greater influence, rewarding active
network participation over passive holding.

\paragraph{Proposal lifecycle.} Proposals follow a three-phase process:
discussion (7 days), voting (7 days), and execution (7-day timelock). A security
council (5-of-9 multi-sig) can veto malicious proposals and pause the network in
emergencies.

\subsection{Fee Reduction}
\label{sec:fees}

Holding HANZO in a fee-reduction contract on platform.hanzo.ai reduces API fees
according to a tiered schedule:

\begin{table}[H]
\centering
\small
\begin{tabular}{rrl}
\toprule
\textbf{HANZO Held} & \textbf{Fee Discount} & \textbf{Tier} \\
\midrule
$\geq$ 1,000 & 10\% & Bronze \\
$\geq$ 10,000 & 20\% & Silver \\
$\geq$ 100,000 & 30\% & Gold \\
$\geq$ 500,000 & 40\% & Platinum \\
$\geq$ 1,000,000 & 50\% & Diamond \\
\bottomrule
\end{tabular}
\caption{Fee discount tiers. Tokens in the fee-reduction contract are liquid (withdrawable with 7-day cooldown) but count toward governance weight.}
\label{tab:fees}
\end{table}

\subsection{Revenue Sharing}
\label{sec:revshare}

HANZO stakers receive a share of platform revenue. Revenue is collected in HANZO
(from compute payments) and distributed weekly:

\begin{equation}
r_i = \frac{s_i}{\sum_j s_j} \cdot R \cdot (1 - \beta - \zeta),
\end{equation}
where $R$ is total weekly revenue, $\beta = 0.10$ is the treasury share, and
$\zeta = 0.25$ is the burn rate. The remaining 65\% flows to stakers.

% ============================================================================
